\vspace{1.2cm}
\phantomsection
\addcontentsline{toc}{chapter}{Liste des abréviations}
{\LARGE\bfseries\color{insatblue} Liste des abréviations\par}\vspace{2pt}\hrule\vspace{10pt}

\begin{small}
\begin{tabular}{@{}p{1.6cm}p{5.3cm}@{\hspace{0.6cm}}p{1.6cm}p{5.3cm}@{}}
\toprule
AE & Auto-encodeur (\emph{auto-encoder}) & LLR & Log-rapport de vraisemblance \\
AWGN & Bruit blanc gaussien additif & MCS & Schéma de modulation et de codage \\
BER & Taux d'erreur binaire & MSE & Erreur quadratique moyenne \\
CE & Entropie croisée (\emph{cross-entropy}) & PSNR & Rapport signal sur bruit crête \\
CIF & \emph{Common Intermediate Format} (352$\times$288) & QAM & Modulation d'amplitude en quadrature \\
CNN & Réseau de neurones convolutif & QPSK & Modulation par déplacement de phase à 4 états \\
CR & Renouvellement conditionnel & ReLU & \emph{Rectified Linear Unit} \\
CSI & Information d'état du canal & RGB & Rouge-vert-bleu \\
DCT & Transformée en cosinus discrète & SNR & Rapport signal sur bruit \\
GOP & Groupe d'images (\emph{group of pictures}) & SSIM & Indice de similarité structurelle \\
JPEG & \emph{Joint Photographic Experts Group} & VGG & \emph{Visual Geometry Group} (architecture) \\
JSCC & Codage conjoint source-canal & ZF & Forçage à zéro (\emph{zero-forcing}) \\
\bottomrule
\end{tabular}
\end{small}

\vspace{14pt}
{\large\bfseries\color{insatblue} Principales notations\par}\vspace{8pt}

\begin{small}
\begin{tabular}{@{}p{2.2cm}p{12.6cm}@{}}
\toprule
$n$ & Nombre de valeurs réelles de l'image source, $n = 3HW$ \\
$k$ & Nombre de symboles complexes émis sur le canal \\
$k/n$ & Rapport de bande ; point de fonctionnement principal $1/12$ \\
$c$ & Nombre de canaux latents effectivement transmis ($k/n = c/96$) \\
$\sigma^2$ & Variance du bruit complexe, $\sigma^2 = 10^{-\mathrm{SNR}/10}$ \\
$h$ & Coefficient de canal de Rayleigh, $h \sim \mathcal{CN}(0,1)$ \\
$x,\ \hat{x}$ & Image émise et image reconstruite au récepteur \\
$z$ & Tenseur latent produit par l'encodeur \\
$\lambda_f,\ \lambda_t$ & Poids des termes de caractéristiques et de tâche dans la perte sémantique \\
\bottomrule
\end{tabular}
\end{small}
